\section{Objectives}

The primary objectives of this work are:

\begin{itemize}
\item To conduct a comprehensive survey of deep learning–based medication recommendation systems, focusing on safety modeling, interpretability, and data requirements.
\item To analyze the limitations of existing approaches, particularly their dependence on dense longitudinal EHR data.
\item To identify research gaps relevant to outpatient and primary-care environments with sparse clinical inputs.
\item To propose a hybrid deep learning architecture that integrates symptom encoding, structured patient features, and drug–drug interaction safety constraints.
\item To design a system that balances predictive accuracy, medication safety, and model interpretability for practical clinical decision support.
\end{itemize}


\section{Proposed Architecture}
We propose a modular system:
\begin{enumerate}
\item \textbf{Preprocessing:} Spell correction, abbreviation expansion, symptom normalization.
\item \textbf{Symptom Encoder:} Lightweight transformer tuned on clinical short-text (or DP-CRNN alternative for low-resource settings).
\item \textbf{Structured Features:} Demographics, comorbid flags, basic vitals if provided.
\item \textbf{Fusion \& Multi-Label Head:} Concatenate embeddings and predict drug classes with sigmoid outputs and label-correlation regularizer.
\item \textbf{Safety Layer:} Rule-based + learned DDI filter (twosides/DrugBank) to remove or down-rank unsafe pairs.
\item \textbf{Explainability:} LIME/SHAP for the classifier + KG/edge trace for safety decisions.
\end{enumerate}
